\section{Endpoint statistics and a charged supercharge: construction hypotheses}
\label{sec:susy-dictionary}
The string interpretation leaves an additional choice open. In an RNS
construction the physical NS/R sectors determine the whole string's spacetime
statistics; Chan--Paton labels alone do not specify independently propagating
spin-$1/2$ endpoint particles. Physical endpoints require boundary fields and
an action. The contact model above chooses two fermionic endpoints, while the
user now also proposes two scalar endpoints for the fermionic string. Either
pair has even fermion parity and can give an odd state with an odd retained
carrier. Mixed scalar/fermion endpoints plus an odd carrier instead have even
total parity. The ordinary SQCD mesino is odd with mixed endpoints and no
additional odd carrier. These are distinct candidates to match, not
interchangeable descriptions already derived from a common theory.

Three charged Dirac leptons have the spin degrees of freedom of six complex
spin-zero fields. The six charged pseudoscalar mesons formed from two up-type
and three down-type hadronizing flavors provide this count. Direct antecedents
exist: Rivero's 2005 composite-boson proposal and 2009 \emph{sBootstrap}
discussion explicitly pursue the three-generation correspondence and the
pointlike-lepton difficulty~\cite{Rivero2005,Rivero2009}. These sources correct
the earlier limited historical search; neither the count nor this research
direction is claimed as new here. The user proposes the pairs
$\mu\leftrightarrow(\pi,K)$, $\tau\leftrightarrow(D,D_s)$ and
$e\leftrightarrow(B,B_c)$, with charge conjugates understood. A dynamical
supersymmetry representation, mass splittings and the top's absence from the
ordinary meson list remain to be explained.

\subsection{A common \texorpdfstring{$B-L$}{B minus L} shift}
For conventional dimensionless number charges, let $X=B-L$. A candidate
supercharge $\mathcal Q_+$ can have the same charge shift in both branches:
\begin{center}
\begin{tabular}{lccc}
\toprule
Transition & Initial $X$ & Final $X$ & Shift\\
\midrule
meson $\to$ antilepton & $0$ & $+1$ & $+1$\\
antidiquark $\to$ quark & $-2/3$ & $+1/3$ & $+1$\\
meson $\to$ lepton & $0$ & $-1$ & $-1$\\
diquark $\to$ antiquark & $+2/3$ & $-1/3$ & $-1$\\
\bottomrule
\end{tabular}
\end{center}
The first two transitions use $\mathcal Q_+$ and the last two its conjugate
$\mathcal Q_-$. The corresponding color representations agree. With suitable
flavors the electric charges can agree as well. Algebraically,
\begin{equation}
 [X,\mathcal Q_\pm]=\pm\mathcal Q_\pm,
 \qquad [\hbar X,\mathcal Q_\pm]=\pm\hbar\mathcal Q_\pm.
\end{equation}
The second expression makes explicit the convention for a dimensionful
generator; the first does not set $\hbar$ to one. If these are spacetime
supercharges, $X$ acts as an R symmetry. Their conjugate charges allow a
neutral translation generator in their anticommutator. This is an admissible
charge ledger, not yet a representation or interaction law.

The two positive shifts require respectively
$(\Delta B,\Delta L)=(0,-1)$ and $(1,0)$.
One fixed charge eigenoperator cannot implement both while carrying definite
charges under independently conserved $B$ and $L$. In a retained-carrier
realization, an exact nonzero continuous $X$ charge also requires distinguishable
carrier/anticarrier sectors; a self-conjugate Majorana carrier cannot carry it.

There is a further, useful constraint on a \emph{linear, exact} $N=1$
interpretation. Choose the convention that a left-chiral scalar of charge $x$
has Weyl partner charge $x-1$. A charged Dirac lepton uses left-handed Weyl
fields with $X=-1,+1$, so the corresponding scalars have $X=0,+2$.
Equivalently the two scalar particles with the lepton's electric charge have
$X=0,-2$. Two ordinary mesons, both $X=0$, do not fill this multiplet with
conserved physical $B-L$. Possible construction choices include an additional
scalar with two units of lepton number, breaking or reassigning the R charge,
or a nonlinear/emergent supersymmetry without these simple observed-state
multiplets. The constraint has these explicit premises; it is not an
exclusion of the broader string proposal.

\subsection{The pion--muon--Goldstino conjecture}
Volkov and Akulov proposed a neutrino as the Goldstone fermion of nonlinear
supersymmetry~\cite{VolkovAkulov}. Their supplied two-page paper does not
identify a pion--muon partner pair. We retain the additional conjecture
\begin{equation}
 \pi^-\longrightarrow\mu^-+\bar G,
 \qquad \bar G\ \hbox{tentatively identified with }\bar\nu_\mu,
\end{equation}
as a broken-supersymmetry scalar-partner transition. The assignments
$X(\mu^-)=-1$ and $X(\bar G)=+1$ conserve $B-L$; the emitted state is
electrically neutral. Nonlinear supersymmetry does not itself require a light
scalar partner, so the proposed pion identification is additional physical
content. Its next substantive test would specify the Goldstino coupling and
compare pion leptonic amplitudes, their helicity dependence and the
muon/electron channel ratio. No such amplitude or decay rate is calculated
in this revision.
